\section{MAMBA-МОДЕЛИ В МАШИННОМ ПЕРЕВОДЕ}

\subsection{Машинный перевод как задача моделирования последовательностей}

Нейронный машинный перевод обычно формулируется как задача условного языкового моделирования.
Пусть исходное предложение задано последовательностью токенов
$x = (x_1, \ldots, x_m)$, а перевод --- последовательностью
$y = (y_1, \ldots, y_n)$.
Тогда модель должна оценивать вероятность перевода
\begin{equation}
    p(y \mid x) = \prod_{t=1}^{n} p(y_t \mid y_{<t}, x).
\end{equation}

На практике для такой задачи часто используют архитектуру encoder-decoder.
Encoder преобразует исходную последовательность в набор скрытых представлений, а decoder по этим представлениям авторегрессионно строит перевод.
В качестве базовой рекуррентной модели в данной работе используется LSTM с механизмом внимания Лионга \cite{luong}.
Такой выбор позволяет сравнить современные архитектуры с классическим подходом к машинному переводу.

Задачу sequence-to-sequence можно сформулировать и без отдельного encoder.
В decoder-only подходе исходное предложение, служебный разделитель и целевой перевод объединяются в одну последовательность.
Модель обучается как языковая модель, то есть предсказывает следующий токен по предыдущим токенам, но функция потерь учитывается только на токенах перевода.
В этом случае исходное предложение играет роль префикса, который задает условие для генерации.
Такая схема используется в современных авторегрессионных языковых моделях и позволяет строить перевод одной моделью без явного разделения на encoder и decoder.

В рекуррентной модели скрытое состояние обновляется последовательно.
В самом общем виде это можно записать так:
\begin{equation}
    h_t = f(h_{t-1}, x_t), \qquad y_t = g(h_t).
\end{equation}
Именно наличие состояния делает LSTM и Mamba близкими по общей идее: обе модели накапливают информацию о прошлых токенах.
Различие заключается в том, как это состояние параметризуется и насколько эффективно его можно вычислять при обучении.

\subsection{Механизм внимания и Transformer}

Модели Transformer стали основой многих современных систем машинного перевода \cite{transformer}.
Их ключевой компонент --- механизм самовнимания, который позволяет каждому токену напрямую учитывать другие токены последовательности.
Для матриц запросов $Q$, ключей $K$ и значений $V$ самовнимание записывается следующим образом:
\begin{equation}
    \operatorname{Attention}(Q, K, V) =
    \operatorname{softmax}\left(\frac{QK^T}{\sqrt{d}}\right)V,
\end{equation}
где $d$ --- размерность векторов ключей.

В encoder-decoder Transformer encoder строит контекстные представления исходного предложения, а decoder использует два вида внимания: причинное самовнимание по уже сгенерированным токенам и внимание к выходам encoder.
Такой подход хорошо работает в машинном переводе, но самовнимание требует попарного сравнения токенов.
Поэтому вычислительная сложность и расход памяти растут квадратично по длине последовательности.
Это ограничение особенно заметно при работе с длинными текстами.

\subsection{Mamba и модели пространства состояний}

Mamba относится к моделям, которые используют скрытое состояние для обработки последовательности \cite{mamba}.
Базовую идею удобно описать через дискретную модель пространства состояний:
\begin{equation}
    \begin{gathered}
        h_t = \bar{A}h_{t-1} + \bar{B}x_t, \\
        y_t = Ch_t + Dx_t,
    \end{gathered}
\end{equation}
где $h_t$ --- скрытое состояние, $x_t$ --- входной токен, $y_t$ --- выход, а матрицы $\bar{A}$, $\bar{B}$, $C$ и $D$ задают обновление состояния и построение выхода.
Если эти параметры фиксированы для всех токенов, модель одинаково обрабатывает все позиции последовательности.
Это удобно для вычислений, но ограничивает способность модели избирательно запоминать важные элементы текста.

Mamba решает эту проблему за счет входозависимых параметров.
Для каждого токена строятся свои значения $\Delta_t$, $B_t$ и $C_t$:
\begin{equation}
    [\delta_t, B_t, C_t] = W_x x_t + b_x, \qquad
    \Delta_t = \operatorname{softplus}(W_{\Delta}\delta_t + b_{\Delta}).
\end{equation}
После этого состояние обновляется по правилу
\begin{equation}
    h_t = \bar{A}_t h_{t-1} + \bar{B}_t x_t, \qquad
    y_t = C_t h_t + D x_t,
\end{equation}
где $\bar{A}_t$ и $\bar{B}_t$ зависят от $\Delta_t$.
Благодаря этому модель может по-разному обрабатывать разные токены: одни сохранять в состоянии, а другие быстро забывать.

Из-за входозависимых параметров Mamba нельзя заранее представить как одну фиксированную свертку.
Тем не менее модель остается быстрой благодаря selective scan --- параллельному сканированию, которое выполняет рекуррентные обновления за линейное время и эффективно использует современные графические ускорители.
Псевдокод блока Mamba приведен в листинге \ref{listing:mamba_pseudocode}.

\begin{lstlisting}[language=Python, caption={Основные операции блока Mamba}, label={listing:mamba_pseudocode}]
Input: x : (B, L, D)
Parameters: A : (D, N)

x, z : (B, L, E * D) <- Linear_in(x)
x : (B, L, E * D) <- SiLU(CausalConv1D(x))

B : (B, L, N) <- s_B(x)
C : (B, L, N) <- s_C(x)
Delta : (B, L, E * D) <- softplus(Parameter + s_Delta(x))

A_bar, B_bar <- discretize(Delta, A, B)
y : (B, L, E * D) <- selective_scan(A_bar, B_bar, C)(x)
y <- y * SiLU(z)

Output: Linear_out(y)
\end{lstlisting}

Для машинного перевода это свойство важно, потому что не все токены одинаково значимы для построения перевода.
Имена собственные, глаголы и числовые значения часто нужно сохранять точнее, чем служебные слова или отдельные знаки препинания.
Поэтому избирательное обновление состояния может быть полезной альтернативой полному самовниманию.

\subsection{Mamba только с декодером}

В классической схеме encoder-decoder исходное предложение и перевод обрабатываются разными частями модели.
В Mamba только с декодером эта граница задается не архитектурой, а форматом последовательности.
Исходное предложение, разделитель и перевод объединяются в один авторегрессионный контекст:
\begin{equation}
    z = (x_1, \ldots, x_m, \operatorname{sep}, y_1, \ldots, y_n).
\end{equation}
Тогда условная вероятность перевода записывается как
\begin{equation}
    p(y \mid x) = \prod_{t=1}^{n}
    p(y_t \mid x_1, \ldots, x_m, \operatorname{sep}, y_{<t}).
\end{equation}

При обучении модель предсказывает следующий токен для всей объединенной последовательности, но ошибка учитывается только на части, соответствующей переводу:
\begin{equation}
    \mathcal{L} = - \sum_{t=1}^{n}
    \log p(y_t \mid x_1, \ldots, x_m, \operatorname{sep}, y_{<t}).
\end{equation}
Во время генерации исходное предложение и разделитель подаются как префикс, после чего модель последовательно добавляет токены перевода.

Такой подход хорошо согласуется с причинной природой Mamba: состояние после обработки исходного предложения становится сжатым представлением условия для дальнейшей генерации.
Преимущество decoder-only схемы состоит в единой авторегрессионной постановке задачи и более простой структуре модели.
Ограничение состоит в том, что связь с исходным предложением проходит через общий последовательный контекст, а не через отдельный слой внимания к выходам encoder.
